\chapter{Software diagrams: method and legend}\label{ch:val-diagrams}

Chapter~\ref{ch:val-functions} ends each symbol's record with a diagram of that symbol. This chapter
says where the diagrams come from, how to read one, and records the verification of the interactive
chart view added on 27 September 2026.

\section{What was parsed}

The page is a single HTML file. Its own script was separated from the inlined chart library and
parsed with a JavaScript parser (acorn, ECMAScript 2022). Every function node in the resulting
syntax tree was collected: function declarations, function expressions and arrow functions, whether
named directly or by the variable, property or assignment they are bound to. 830 named functions
were found, 614 of them at the top level of the script; the remainder are defined inside another
function and are named on the diagram of the function that contains them. Every name declared at the
top level that is not a function was treated as shared state; 124 such names were found.

The diagrams are generated from the source by parsing it, not written by hand, so they describe the
code as it is rather than as it was intended.

\section{How the arrows were derived}

For each function, every call expression in its own body was collected, excluding calls that belong
to a nested \emph{named} function, which gets its own entry. A call written inside an anonymous
callback counts as a call by the function that contains it, because that is how the code reads. An
arrow is drawn only between symbols defined in the page; calls to the browser, to the standard
library and to array methods are not drawn, as they would swamp the diagram. A function that calls
itself is not drawn as calling itself.

\section{How to read a diagram}

Left: the functions in the page that call this one. Nothing on the left means the function is either
an entry point --- an event handler, the boot sequence --- or is reached through a table of
definitions rather than by name. Centre: the symbol itself, with its arguments as they are written
in the source, whether it returns a value, and the named helpers defined inside it. Right: the
functions it calls. Below: in blue, the shared state it mentions; in amber, the effects detected in
its body --- writing to the document, producing a download, reading a file, browser storage,
hashing, drawing an interactive figure, wiring events, notifying the user. A symbol with no chips
computes from its arguments and changes nothing outside itself.

For a declared value rather than a function, the left column lists the functions that mention it.
At most seven boxes are drawn on each side; any remainder is counted in a final box, and the true
degree is stated under the centre box.

\section{What a diagram does not tell you}

The arrows are static: they show that a call appears in the source, not how often it runs or whether
a branch is taken. A function reached only through an event handler or through a table of
definitions --- the graph renderers, for instance --- shows no callers even though it runs.

\section{Validation record: the interactive chart view}

An on-screen interactive view of the charts was added on 27 September 2026, drawn with an inlined
chart library so the page remains a single offline file. \texttt{svgPlot} remains the single source
of truth: it still draws every chart, and its SVG is what the \texttt{.svg} download, the
\texttt{.png} rasteriser, the ZIP and the export atlas take. The interactive figure is rebuilt from
the same options object \texttt{svgPlot} was handed. The drawn SVG is not discarded when a figure is
shown: it is moved into a hidden holder that stays first in document order, so an export that reads
the chart out of the page still reads the drawn SVG.

Both the page as published and the page proposed to replace it were loaded in a browser, given the
same 55 demonstration runs, and asked for every entry in the download catalogue. Each entry was
hashed and the two sets compared.

\begin{center}\small
\begin{table}[H]\centering\small
\begin{tabular}{|l|r|}
\hline
Page validated & \texttt{index.html}, 2\,021\,629\,B \\ \hline
Its SHA-256 & \texttt{327183b2\ldots cedd3753} \\ \hline
Compared against & the page then published, 940\,696\,B, \texttt{8df4472e\ldots b33c770b} \\ \hline
Runs read from the demonstration set & 55 (\texttt{LC-DEMO-01}) \\ \hline
Catalogue entries produced by both pages & 101 \\ \hline
Export entries written by both pages & 139 \\ \hline
Byte-identical between the two pages & \vPASS{} 138 of 139 \\ \hline
The one difference & \texttt{README.txt}: its own \texttt{Built:} time, 1.3\,s apart \\ \hline
Entries that failed to draw & 0 \\ \hline
Graphs drawn through the interactive layer & 13 of 25 \\ \hline
Graphs kept on their drawn SVG & 12 of 25 \\ \hline
Control charts drawn interactively & 20 --- every chart with data in this set \\ \hline
Control charts kept on their drawn SVG & 38 of 58 defined \\ \hline
Page errors / console errors & 0 / 0 \\ \hline
Time to catalogue ready, published page / this page & 215\,ms / 240\,ms \\ \hline
\end{tabular}
\caption{Symbols of the area \emph{Validation record: the interactive chart view}}
\end{table}
\end{center}

\begin{notebox}[title={\textbf{This record was re-measured on 28 September 2026}}]
The figures above replace an earlier record taken against a build of 2\,007\,034 bytes. That build
was made from the page as it stood on 25 September and did not carry three fixes that were already
in the published page, so it was never released; the page validated here is the interactive layer
re-applied to the published page itself. Both pages were driven headless with the same 55
demonstration runs, every catalogue item was selected, and the export set each page wrote was hashed
entry by entry. The counts for the interactive layer are measured through the catalogue's own export
path with default options and one instrument, which is why they differ from the earlier record's
23 of 25: that figure was taken in the interface with a target chosen for each graph.
\end{notebox}


\subsection*{Defects found during the work}

\textbf{A page element shadowed the module loader.} The page contains an element with
\texttt{id="exports"}. A browser exposes every element id as a window property, so the chart
library's loader read that element as a module environment and attached itself to it instead of to
the window. The library was then undefined, with no error raised anywhere. The library is now
evaluated with \texttt{exports}, \texttt{module} and \texttt{define} shadowed.

\textbf{A stylesheet rule blanked a correctly built figure.} The rules that style a drawn chart apply
to every SVG inside the chart box, including the layered SVGs an interactive figure is made of,
giving one of them an opaque white background. The figure was built correctly --- the nodes were all
present and countable --- but painted over. Only a screenshot revealed it. The rules are now scoped
away from the interactive figures.

\textbf{An asynchronous draw outlived its own figure.} Drawing finishes asynchronously. Replacing a
figure before it settled left the remainder of that work with nothing to act on, which surfaced as a
page error. Figures are now released before their nodes are replaced, and the late rejection is
handled.

\subsection*{What was deliberately not changed}

The plate heat map and the run acceptance grid draw their own grid of rectangles rather than going
through \texttt{svgPlot}. They keep their drawn SVG, so their exported bytes are untouched. No stored
value, no call, no chart mathematics and no exported byte changed.
